\clearpage\fancyhead[L]{\small COMET.}\hypertarget{native-407-part-lower}{}\label{native-407-part-lower}
\begin{center}\large COMET.\end{center}
\begin{center}\hrule\vspace{4pt}\fontsize{10}{12}\selectfont V. \textsc{Lindenau} and \textsc{Bohnenberger} Zeitschrift für Astronomie u. verwandte Wiss.\par Vol. V. pp. 276..277. March and April 1818.\par\vspace{4pt}\hrule\vspace{4pt}\end{center}
\par\indent I recently derived the following parabolic elements from Dr. \textsc{Olbers}’s observations of the comet from 3 to 28 March. They differ only slightly from those \textsc{Olbers} himself obtained from his observations of 3 and 13 March and the observation made in Marseille on 18 January, which is still unknown to me.

\begin{center}\fontsize{10}{12}\selectfont
Passage through perihelion, 26,8991 February 1818, Göttingen time.
\par\smallskip\begin{tabular}{lrr}
Logarithm of the perihelion distance & $\cdot\;\cdot\;\cdot\;\cdot\;\cdot\;\cdot\;\cdot\;\cdot$ & $0.07937$\\
Inclination of the orbit & $\cdot\;\cdot\;\cdot\;\cdot\;\cdot\;\cdot\;\cdot\;\cdot\;\cdot\;\cdot\;\cdot\;\cdot\;\cdot\;\cdot\;\cdot$ & $90^\circ\;0'\;0''$\\
Ascending node & $\cdot\;\cdot\;\cdot\;\cdot\;\cdot\;\cdot\;\cdot\;\cdot\;\cdot\;\cdot\;\cdot\;\cdot\;\cdot\;\cdot$ & $70\;5\;12$\\
\end{tabular}
\par\smallskip Longitude of perihelion, taking the motion as prograde $183\;22\;58$
\end{center}
\SourceLeaf{408}\par\indent The orbit’s perpendicular orientation to the ecliptic followed naturally from the observations, but it will probably be modified somewhat as observations continue. Only then will it become clear whether the motion is prograde or retrograde with respect to the ecliptic. To make observations easier, I have calculated the following ephemeris from these elements. The time is $14^{\mathrm u}20^{\mathrm m}$ in Göttingen. The brightness column assumes borrowed light, $=1$ at distance 1 from Earth and Sun.

\begin{center}\fontsize{10}{12}\selectfont\setlength{\tabcolsep}{3pt}
\begin{tabular}{ll|ll|lll|l}
\multicolumn{2}{c|}{1818} & \multicolumn{2}{c|}{R.A.} & \multicolumn{3}{c|}{Decl.} & Brightness \\
\hline
April & 6  & $298^\circ$ & $58'$ & $12^\circ$ & $57'$ & N. & 0.396 \\
      & 10 & 297 & 50 & 11 & 2  &    & 0.418 \\
      & 14 & 296 & 27 & 8  & 54 &    & 0.443 \\
      & 18 & 294 & 47 & 6  & 30 &    & 0.471 \\
      & 22 & 292 & 44 & 3  & 48 &    & 0.503 \\
      & 26 & 290 & 17 & 0  & 43 &    & 0.535 \\
      & 30 & 287 & 21 & 2  & 46 & S. & 0.568 \\
May   & 4  & 283 & 53 & 6  & 41 &    & 0.596 \\
      & 8  & 279 & 47 & 10 & 58 &    & 0.616 \\
      & 12 & 275 & 4  & 15 & 33 &    & 0.621 \\
\end{tabular}
\end{center}

\par\indent Moonlight will probably prevent any further observations after that. By the time the moonlight is gone, the comet will have moved too far south, for us at least.

\par\indent I myself have so far observed the comet only once and compared it with a star of the 9\textsuperscript{th} magnitude that is absent from the Hist. cél. I leave the determination of its position for later. Its position will be approximately $300^\circ\,3'$ R.A. and $15^\circ\,9'$ declination.

\begin{center}\fontsize{9}{11}\selectfont\setlength{\tabcolsep}{3pt}\begin{tabular}{lrrrll}
1 April. & $15^{\mathrm u}$ & $45^{\mathrm m}$ & $28^{\mathrm s}$ & Sidereal time, & \shortstack[l]{the comet follows the star by\\$14''98$ in time}\\
& $15$ & $44$ & $56$ & — & \shortstack[l]{the comet is $4'\,22''$\\farther south than the star.}\\
\end{tabular}\end{center}

\par\indent The declination is probably considerably less accurate than the right ascension.

\Needspace{8\baselineskip}\par\indent [Manuscript:] Observations of this star
\begin{center}\fontsize{10}{12}\selectfont\begin{tabular}{lrl}
August 1819 & 27 & $20^{\mathrm u}\;0^{\mathrm m}\;11^{\mathrm s}49$\\
Sept. & 4 & $11.48$\\
& 7 & $11.38$\\
& 24 & $11.30$\\
\end{tabular}\end{center}

